\section{Predykcja}\label{sec:predykcja}

Wynik automatycznej ekstrakcji dla każdego schematyzmu przyjmuje postać sekwencji rekordów o identycznej strukturze pól co dane referencyjne.
Ponieważ model ekstrakcji zwraca wartości w oryginalnej formie źródłowej (łacińskiej, rosyjskiej bądź niemieckiej), a baza referencyjna projektu głównego posługuje się polskimi odpowiednikami i skrótami, konieczny jest dodatkowy etap tłumaczenia, w dalszym ciągu określany jako \textit{parsowanie}.
Etap ten opiera się na trzech słownikach mapujących formy źródłowe na znormalizowane wartości polskie: dla materiału budowlanego (np.\ \textit{mur.}~$\rightarrow$~\textit{mr}, \textit{lign.}~$\rightarrow$~\textit{dr}), dla wezwań kościelnych (np.\ \textit{Assumptionis B.M.V.}~$\rightarrow$~\textit{Wniebowzięcie NMP}) oraz dla nazw dekanatów (np.\ \textit{Decanatus Gedanensis}~$\rightarrow$~\textit{Gdańsk}).
Aby uwzględnić warianty pisowni i drobne błędy w odczycie, dopasowanie w słownikach odbywa się z wykorzystaniem przybliżonego porównania łańcuchów z konfigurowalnym progiem tolerancji.
W rezultacie ewaluacja prowadzona jest dwutorowo: raz na poziomie formy źródłowej (\textit{source}), a raz na poziomie formy przetłumaczonej (\textit{parsed}), co pozwala oddzielnie mierzyć jakość samej ekstrakcji od jakości tłumaczenia na język docelowy bazy.
Porównanie obu sekwencji, referencyjnej i wygenerowanej, nie jest jednak trywialne: mogą się one różnić długością (model mógł pominąć pewne parafie lub wygenerować wpisy nieistniejące w źródle), a kolejność rekordów nie musi się pokrywać.
Problem ten rozwiązano za pomocą algorytmu węgierskiego, który znajduje optymalne przyporządkowanie parami między rekordami referencyjnymi a przewidywanymi, minimalizując łączną niepodobność.
Podobieństwo dwóch rekordów obliczane jest jako średnia ważona podobieństw poszczególnych pól, gdzie każde pole (dekanat, parafia, wezwanie, materiał budowlany) ma przypisaną wagę odzwierciedlającą jego znaczenie.
Podobieństwo pojedynczego pola mierzone jest znormalizowaną odległością Levenshteina po uprzednim sprowadzeniu obu wartości do formy porównywalnej: usunięciu znaków diakrytycznych, ujednoliceniu wielkości liter i normalizacji białych znaków.
Rekordy, których wzajemne podobieństwo nie przekracza ustalonego progu, nie są łączone w pary i trafiają odpowiednio do kategorii wpisów pominiętych (po stronie referencji) lub nadmiarowych (po stronie predykcji).

